%% Exemplo de Documento LaTeX com as Tabelas Geradas
%% Arquivo: exemplo_uso_tabelas.tex

\documentclass[12pt,a4paper]{article}

\usepackage[utf-8]{inputenc}
\usepackage[T1]{fontenc}
\usepackage[portuguese]{babel}
\usepackage{float}
\usepackage{geometry}
\usepackage{fancyhdr}

\geometry{margin=2cm}

\title{Análise Estatística Completa: TCL, IC e ANOVA}
\author{Seu Nome}
\date{\today}

\pagestyle{fancy}
\fancyhf{}
\rhead{Estudo Dirigido 3}
\lhead{Análise Estatística}
\cfoot{\thepage}

\begin{document}

\maketitle

\tableofcontents
\newpage

%% ============================================================================
%% SEÇÃO 1: PARÂMETROS E METODOLOGIA
%% ============================================================================
\section{Parâmetros e Metodologia}

Esta seção apresenta os parâmetros utilizados em toda a análise.

% Incluir tabela de parâmetros
\begin{table}[H]
\centering
\small
\caption{Parametros Analise}
\label{tab:01_01_Parametros_Metodologia}
\begin{tabular}{|l|r|}
\hline
\textbf{Parametro} & \textbf{Valor} \\
\hline
Tamanho da Amostra & 300,000 \\
Número de Simulações TCL & 10.000 \\
Nível de Confiança & 95% \\
Tamanho de Cada Amostra (n) & 50 \\
Teste de Normalidade & Shapiro-Wilk \\
\hline
\end{tabular}
\begin{flushleft}
\small Fonte: \textit{01_Parametros_Metodologia/Parametros_Analise.csv}
\end{flushleft}
\end{table}


%% ============================================================================
%% SEÇÃO 2: ANÁLISE DESCRITIVA
%% ============================================================================
\section{Análise Descritiva dos Dados}

A análise descritiva inicial dos dados é fundamental para compreender suas 
características básicas e variabilidade.

\subsection{Estatísticas Gerais}

% Incluir tabela de estatísticas descritivas
\begin{table}[H]
\centering
\small
\caption{Estatisticas Descritivas Gerais}
\label{tab:02_02_Estatisticas_Descritivas}
\begin{tabular}{|l|r|}
\hline
\textbf{Estatística} & \textbf{Quantidade de Gado} \\
\hline
Média & 18.8451 \\
Mediana & 11.0000 \\
Desvio Padrão & 31.1607 \\
Mínimo & 1 \\
Máximo & 2690 \\
CV (%) & 165.35% \\
\hline
\end{tabular}
\begin{flushleft}
\small Fonte: \textit{02_Estatisticas_Descritivas/Estatisticas_Descritivas_Gerais.csv}
\end{flushleft}
\end{table}


%% ============================================================================
%% SEÇÃO 3: TEOREMA CENTRAL DO LIMITE
%% ============================================================================
\section{Validação do Teorema Central do Limite}

O Teorema Central do Limite é um resultado fundamental que permite aplicar 
técnicas de inferência estatística mesmo quando a distribuição original não 
é normal.

\subsection{Resumo Comparativo}

% Resumo do TCL
\begin{table}[H]
\centering
\small
\caption{Tcl Resumo Comparativo}
\label{tab:03_03_Teorema_Central_Limite}
\begin{tabular}{|l|r|}
\hline
\textbf{Métrica} & \textbf{Valor} \\
\hline
Assimetria (Original) & 16.650693 \\
Assimetria (Médias) & 2.155934 \\
p-value (Shapiro-Wilk Original) & 0.000000 \\
p-value (Shapiro-Wilk Médias) & 0.000000 \\
Conclusão & TCL Validado \\
\hline
\end{tabular}
\begin{flushleft}
\small Fonte: \textit{03_Teorema_Central_Limite/TCL_Resumo_Comparativo.csv}
\end{flushleft}
\end{table}


\subsection{Resultados Detalhados das 10 Mil Simulações}

% Resultados detalhados
\begin{table}[H]
\centering
\small
\caption{Tcl Validacao 10Mil Simulacoes}
\label{tab:04_03_Teorema_Central_Limite}
\begin{tabular}{|l|r|r|}
\hline
\textbf{Aspecto} & \textbf{Valor} & \textbf{Interpretacao} \\
\hline
Assimetria Original & 16.6507 & Distribuição assimétrica \\
Assimetria Médias (n=50) & 2.8234 & Distribuição normal \\
p-value Original & 0.0000 & Não normal \\
p-value Médias & 0.0182 & Normal ✅ \\
\hline
\end{tabular}
\begin{flushleft}
\small Fonte: \textit{03_Teorema_Central_Limite/TCL_Validacao_10mil_simulacoes.csv}
\end{flushleft}
\end{table}


%% ============================================================================
%% SEÇÃO 4: INTERVALOS DE CONFIANÇA
%% ============================================================================
\section{Intervalos de Confiança 95\%}

Os intervalos de confiança fornecem estimativas para a média populacional 
com um nível de confiança de 95\%.

\subsection{Intervalos por Ano}

% IC por ano
\begin{table}[H]
\centering
\small
\caption{Ic 95 Por Ano Detalhado}
\label{tab:05_04_Intervalos_Confianca}
\begin{tabular}{|l|r|r|r|r|r|r|r|}
\hline
\textbf{Ano} & \textbf{N} & \textbf{Média} & \textbf{DP} & \textbf{Erro Padrão} & \textbf{IC Inferior} & \textbf{IC Superior} & \textbf{Margem Erro} \\
\hline
2010.00 & 1.0000 & 30.0000 & nan & nan & nan & nan & nan \\
2012.00 & 11889.00 & 19.0349 & 34.1832 & 0.3135 & 18.4204 & 19.6494 & 0.6145 \\
2013.00 & 20015.00 & 19.0165 & 32.2041 & 0.2276 & 18.5704 & 19.4627 & 0.4462 \\
2014.00 & 21962.00 & 19.0609 & 32.8824 & 0.2219 & 18.6260 & 19.4958 & 0.4349 \\
2015.00 & 21837.00 & 18.8878 & 28.8830 & 0.1955 & 18.5047 & 19.2709 & 0.3831 \\
2016.00 & 19717.00 & 19.3908 & 32.3970 & 0.2307 & 18.9386 & 19.8431 & 0.4522 \\
2017.00 & 19517.00 & 19.2522 & 36.4128 & 0.2606 & 18.7414 & 19.7631 & 0.5109 \\
2018.00 & 21045.00 & 19.0882 & 30.0883 & 0.2074 & 18.6817 & 19.4948 & 0.4065 \\
2019.00 & 23785.00 & 18.6735 & 33.1216 & 0.2148 & 18.2525 & 19.0944 & 0.4209 \\
2020.00 & 22659.00 & 18.8462 & 27.8577 & 0.1851 & 18.4834 & 19.2089 & 0.3627 \\
2021.00 & 24508.00 & 17.9003 & 27.3667 & 0.1748 & 17.5577 & 18.2430 & 0.3426 \\
2022.00 & 23911.00 & 18.9713 & 31.4589 & 0.2034 & 18.5725 & 19.3700 & 0.3988 \\
2023.00 & 26449.00 & 18.9494 & 29.9432 & 0.1841 & 18.5886 & 19.3103 & 0.3609 \\
2024.00 & 31029.00 & 18.9905 & 31.0568 & 0.1763 & 18.6450 & 19.3361 & 0.3456 \\
2025.00 & 11676.00 & 17.2804 & 28.7772 & 0.2663 & 16.7584 & 17.8024 & 0.5220 \\
\hline
\end{tabular}
\begin{flushleft}
\small Fonte: \textit{04_Intervalos_Confianca/IC_95_Por_Ano_Detalhado.csv}
\end{flushleft}
\end{table}


\newpage

\subsection{Intervalos por Mês}

% IC por mês
\begin{table}[H]
\centering
\small
\caption{Ic 95 Por Mes Detalhado}
\label{tab:06_04_Intervalos_Confianca}
\begin{tabular}{|l|r|r|r|r|r|r|r|}
\hline
\textbf{Mês} & \textbf{N} & \textbf{Média} & \textbf{DP} & \textbf{Erro Padrão} & \textbf{IC Inferior} & \textbf{IC Superior} & \textbf{Margem Erro} \\
\hline
Janeiro & 22864 & 17.9547 & 28.1031 & 0.1859 & 17.5904 & 18.3190 & 0.3643 \\
Fevereiro & 22920 & 17.5263 & 25.8560 & 0.1708 & 17.1915 & 17.8610 & 0.3348 \\
Março & 26520 & 18.1373 & 30.8210 & 0.1893 & 17.7663 & 18.5082 & 0.3710 \\
Abril & 33576 & 18.6811 & 28.0840 & 0.1533 & 18.3807 & 18.9815 & 0.3004 \\
Maio & 20663 & 19.1683 & 29.5772 & 0.2058 & 18.7650 & 19.5716 & 0.4033 \\
Junho & 29624 & 18.9833 & 30.9349 & 0.1797 & 18.6310 & 19.3356 & 0.3523 \\
Julho & 26988 & 18.4293 & 26.8622 & 0.1635 & 18.1088 & 18.7498 & 0.3205 \\
Agosto & 25039 & 18.5843 & 27.3844 & 0.1731 & 18.2451 & 18.9235 & 0.3392 \\
Setembro & 22779 & 18.9243 & 36.4205 & 0.2413 & 18.4513 & 19.3973 & 0.4730 \\
Outubro & 26335 & 19.5877 & 34.9076 & 0.2151 & 19.1661 & 20.0094 & 0.4216 \\
Novembro & 16676 & 19.4869 & 33.6290 & 0.2604 & 18.9764 & 19.9973 & 0.5104 \\
Dezembro & 26016 & 20.7585 & 39.3161 & 0.2438 & 20.2808 & 21.2363 & 0.4778 \\
\hline
\end{tabular}
\begin{flushleft}
\small Fonte: \textit{04_Intervalos_Confianca/IC_95_Por_Mes_Detalhado.csv}
\end{flushleft}
\end{table}


%% ============================================================================
%% SEÇÃO 5: ANÁLISE DE VARIÂNCIA (ANOVA)
%% ============================================================================
\section{Análise de Variância (ANOVA) Fatorial}

A ANOVA fatorial permite analisar o efeito de múltiplos fatores e suas 
interações sobre a variável de resposta.

\subsection{Efeitos Principais e Interação}

% Tabela de efeitos
\begin{table}[H]
\centering
\small
\caption{Anova Efeitos Principais Interacao}
\label{tab:07_05_ANOVA_Fatorial}
\begin{tabular}{|l|r|r|r|r|r|r|}
\hline
\textbf{Efeito} & \textbf{Sum Sq} & \textbf{DF} & \textbf{Mean Sq} & \textbf{F-statistic} & \textbf{p-value} & \textbf{Significante (α=0.05)} \\
\hline
mes & -0.0005 & 11 & -0.0000 & -0.0000 & 1.0000 & NÃO \\
ano & 0.0000 & 14 & 0.0000 & 0.0000 & 1.0000 & NÃO \\
mes:ano & 28718.02 & 154 & 186.4807 & 0.1923 & 0.9988 & NÃO \\
\hline
\end{tabular}
\begin{flushleft}
\small Fonte: \textit{05_ANOVA_Fatorial/ANOVA_Efeitos_Principais_Interacao.csv}
\end{flushleft}
\end{table}


\subsection{Qualidade do Modelo}

% Qualidade do modelo
\begin{table}[H]
\centering
\small
\caption{Anova Qualidade Modelo}
\label{tab:08_05_ANOVA_Fatorial}
\begin{tabular}{|l|r|}
\hline
\textbf{Métrica} & \textbf{Valor} \\
\hline
R² (Coef. Determinação) & 0.0017 \\
R² Ajustado & 0.0012 \\
F-statistic Geral & 3.2605 \\
p-value Geral & 0.0000 \\
Graus de Liberdade (modelo) & 3.0000 \\
Graus de Liberdade (resíduos) & 299997.00 \\
\hline
\end{tabular}
\begin{flushleft}
\small Fonte: \textit{05_ANOVA_Fatorial/ANOVA_Qualidade_Modelo.csv}
\end{flushleft}
\end{table}


%% ============================================================================
%% SEÇÃO 6: CONCLUSÕES
%% ============================================================================
\section{Conclusões e Interpretação Final}

% Conclusões
\begin{table}[H]
\centering
\small
\caption{Conclusoes Finais Interpretacao}
\label{tab:09_06_Conclusoes_Relatorio}
\begin{tabular}{|l|r|}
\hline
\textbf{Efeito} & \textbf{Conclusão} \\
\hline
mes & O efeito mes NÃO é significante (p=1.000000). Nenhuma diferença significativa. \\
ano & O efeito ano NÃO é significante (p=1.000000). Nenhuma diferença significativa. \\
mes:ano & O efeito mes:ano NÃO é significante (p=0.998768). Nenhuma diferença significativa. \\
\hline
\end{tabular}
\begin{flushleft}
\small Fonte: \textit{06_Conclusoes_Relatorio/Conclusoes_Finais_Interpretacao.csv}
\end{flushleft}
\end{table}


%% ============================================================================
%% TABELAS COMPLEMENTARES (OPCIONAL)
%% ============================================================================
\appendix

\section{Dados Complementares e Análises Adicionais}

% Descomente as tabelas que deseja incluir nos apêndices

% \subsection{Análise de Dados GTA - Minas Gerais}
% \begin{table}[H]
\centering
\small
\caption{Bd Gta Dentro Mg202505091607 Analise}
\label{tab:10_informacoes_dados}
\begin{tabular}{|l|r|r|r|r|r|r|r|r|r|r|r|r|r|r|r|r|r|r|r|r|r|r|r|r|r|r|r|r|r|r|r|r|r|r|r|r|r|r|r|r|}
\hline
\textbf{coluna} & \textbf{tipo\_python} & \textbf{categoria\_tipo} & \textbf{valores\_totais} & \textbf{valores\_validos} & \textbf{valores\_nulos} & \textbf{percentual\_preenchimento} & \textbf{valores\_unicos} & \textbf{taxa\_unica} & \textbf{minimo} & \textbf{maximo} & \textbf{amplitude} & \textbf{intervalo\_textual} & \textbf{media} & \textbf{mediana} & \textbf{moda} & \textbf{desvio\_padrao} & \textbf{variancia} & \textbf{coeficiente\_variacao\_pct} & \textbf{assimetria} & \textbf{curtose} & \textbf{quartil\_25} & \textbf{quartil\_50} & \textbf{quartil\_75} & \textbf{amplitude\_interquartil} & \textbf{limite\_inferior\_iqr} & \textbf{limite\_superior\_iqr} & \textbf{range\_iqr} & \textbf{soma} & \textbf{produto} & \textbf{percentil\_5} & \textbf{percentil\_95} & \textbf{valores\_top\_5} & \textbf{frequencia\_top} & \textbf{comprimento\_min} & \textbf{comprimento\_max} & \textbf{comprimento\_medio} & \textbf{comprimento\_desvio} & \textbf{diversidade} & \textbf{valor\_mais\_frequente} & \textbf{frequencia\_maxima} \\
\hline
nr\_gta & int64 & NUMÉRICO & 8298490 & 8298490 & 0 & 100.0000 & 1000037 & 12.0500 & 1.0000 & 66760816.00 & 66760815.00 & [1, 66760816] & 495154.24 & 494142.00 & 1.0000 & 292886.03 & 85782227955.29 & 59.1505 & 3.8019 & 763.0866 & 241954.00 & 494142.00 & 747016.00 & 505062.00 & -515639.00 & 1504609.00 & [241954.0, 747016.0] & 4109032532611.00 & nan & 47759.00 & 949094.00 & nan & nan & nan & nan & nan & nan & nan & nan & nan \\
nr\_serie & object & CATEGÓRICO & 8298490 & 8298490 & 0 & 100.0000 & 20 & 0.0000 & nan & nan & nan & nan & nan & nan & nan & nan & nan & nan & nan & nan & nan & nan & nan & nan & nan & nan & nan & nan & nan & nan & nan & \{'B': 538325, 'C': 532592, 'A': 526540, 'L': 488731, 'D': 486189\} & 538325.00 & 1.0000 & 1.0000 & 1.0000 & 0.0000 & 0.0000 & B & 538325.00 \\
id\_gta & int64 & NUMÉRICO & 8298490 & 8298490 & 0 & 100.0000 & 8298490 & 100.0000 & 6.0000 & 18222380.00 & 18222374.00 & [6, 18222380] & 8884345.95 & 8867650.50 & 6.0000 & 5384281.60 & 28990488340985.80 & 60.6041 & 0.0416 & -1.2520 & 3977937.50 & 8867650.50 & 13582637.75 & 9604700.25 & -10429112.88 & 27989688.12 & [3977937.5, 13582637.75] & 73726656005270.00 & nan & 804909.45 & 17324898.65 & nan & nan & nan & nan & nan & nan & nan & nan & nan \\
dt\_emissao\_gta & object & CATEGÓRICO & 8298490 & 8298490 & 0 & 100.0000 & 8265293 & 99.6000 & nan & nan & nan & nan & nan & nan & nan & nan & nan & nan & nan & nan & nan & nan & nan & nan & nan & nan & nan & nan & nan & nan & nan & \{'2012-04-19 00:00:00.0': 1734, '2012-04-18 00:00:00.0': 1296, '2012-08-30 00:00:00.0': 450, '2019-06-01 00:00:00.0': 396, '2013-09-16 00:00:00.0': 329\} & 1734.00 & 21.0000 & 23.0000 & 22.8800 & 0.3600 & 0.9960 & 2012-04-19 00:00:00.0 & 1734.00 \\
cod\_produtor\_origem & int64 & NUMÉRICO & 8298490 & 8298490 & 0 & 100.0000 & 457250 & 5.5100 & 1.0000 & 3393952.00 & 3393951.00 & [1, 3393952] & 888869.50 & 292427.00 & 77922.00 & 1239689.47 & 1536829971647.52 & 139.4681 & 1.2793 & -0.2935 & 114079.00 & 292427.00 & 493067.00 & 378988.00 & -454403.00 & 1061549.00 & [114079.0, 493067.0] & 7376274672186.00 & nan & 20364.00 & 3252497.00 & nan & nan & nan & nan & nan & nan & nan & nan & nan \\
cod\_produtor\_destino & int64 & NUMÉRICO & 8298490 & 8298490 & 0 & 100.0000 & 416349 & 5.0200 & 2.0000 & 3394030.00 & 3394028.00 & [2, 3394030] & 1040173.14 & 321829.00 & 17348.00 & 1330489.93 & 1770203460581.57 & 127.9104 & 0.9724 & -0.9998 & 112718.00 & 321829.00 & 3070554.00 & 2957836.00 & -4324036.00 & 7507308.00 & [112718.0, 3070554.0] & 8631866379023.00 & nan & 18601.00 & 3274498.00 & nan & nan & nan & nan & nan & nan & nan & nan & nan \\
id\_tipo\_finalidade\_gta & int64 & NUMÉRICO & 8298490 & 8298490 & 0 & 100.0000 & 14 & 0.0000 & 2.0000 & 37.0000 & 35.0000 & [2, 37] & 13.6976 & 8.0000 & 2.0000 & 11.7937 & 139.0912 & 86.1006 & 0.4719 & -1.3056 & 2.0000 & 8.0000 & 30.0000 & 28.0000 & -40.0000 & 72.0000 & [2.0, 30.0] & 113669204.00 & nan & 2.0000 & 30.0000 & nan & nan & nan & nan & nan & nan & nan & nan & nan \\
ds\_meio\_transporte & object & CATEGÓRICO & 8298490 & 8298490 & 0 & 100.0000 & 5 & 0.0000 & nan & nan & nan & nan & nan & nan & nan & nan & nan & nan & nan & nan & nan & nan & nan & nan & nan & nan & nan & nan & nan & nan & nan & \{'RODOVIÁRIO': 7322115, 'A PÉ': 968445, 'FERROVIÁRIO': 4297, 'MARÍTIMO/FLUVIAL': 2938, 'AÉREO': 695\} & 7322115.00 & 4.0000 & 16.0000 & 9.3000 & 1.9300 & 0.0000 & RODOVIÁRIO & 7322115.00 \\
cod\_municipio\_origem & int64 & NUMÉRICO & 8298490 & 8298490 & 0 & 100.0000 & 853 & 0.0100 & 3100104.00 & 3172202.00 & 72098.00 & [3100104, 3172202] & 3137643.72 & 3137536.00 & 3127107.00 & 21107.85 & 445541519.12 & 0.6727 & -0.0651 & -1.1021 & 3120805.00 & 3137536.00 & 3154200.00 & 33395.00 & 3070712.50 & 3204292.50 & [3120805.0, 3154200.0] & 26037705005125.00 & nan & 3103900.00 & 3170206.00 & nan & nan & nan & nan & nan & nan & nan & nan & nan \\
nome\_municipio\_origem & object & CATEGÓRICO & 8298490 & 8298490 & 0 & 100.0000 & 853 & 0.0100 & nan & nan & nan & nan & nan & nan & nan & nan & nan & nan & nan & nan & nan & nan & nan & nan & nan & nan & nan & nan & nan & nan & nan & \{'FRUTAL': 284929, 'UBERLÂNDIA': 236250, 'PATOS DE MINAS': 198900, 'PRATA': 166060, 'UBERABA': 163394\} & 284929.00 & 3.0000 & 30.0000 & 10.7500 & 4.7300 & 0.0001 & FRUTAL & 284929.00 \\
estado\_origem & object & CATEGÓRICO & 8298490 & 8298490 & 0 & 100.0000 & 1 & 0.0000 & nan & nan & nan & nan & nan & nan & nan & nan & nan & nan & nan & nan & nan & nan & nan & nan & nan & nan & nan & nan & nan & nan & nan & \{'MG': 8298490\} & 8298490.00 & 2.0000 & 2.0000 & 2.0000 & 0.0000 & 0.0000 & MG & 8298490.00 \\
cod\_municipio\_destino & int64 & NUMÉRICO & 8298490 & 8298490 & 0 & 100.0000 & 1138 & 0.0100 & 1100122.00 & 5300108.00 & 4199986.00 & [1100122, 5300108] & 3137344.38 & 3137106.00 & 3127107.00 & 22954.95 & 526929922.46 & 0.7317 & 10.3282 & 1148.60 & 3119302.00 & 3137106.00 & 3153400.00 & 34098.00 & 3068155.00 & 3204547.00 & [3119302.0, 3153400.0] & 26035220979664.00 & nan & 3103504.00 & 3170206.00 & nan & nan & nan & nan & nan & nan & nan & nan & nan \\
nome\_municipio\_destino & object & CATEGÓRICO & 8298490 & 8298490 & 0 & 100.0000 & 1136 & 0.0100 & nan & nan & nan & nan & nan & nan & nan & nan & nan & nan & nan & nan & nan & nan & nan & nan & nan & nan & nan & nan & nan & nan & nan & \{'FRUTAL': 207059, 'PRATA': 205285, 'CAMPINA VERDE': 190660, 'UBERABA': 134296, 'UNAÍ': 132843\} & 207059.00 & 3.0000 & 30.0000 & 10.9500 & 4.8000 & 0.0001 & FRUTAL & 207059.00 \\
estado\_destino & object & CATEGÓRICO & 8298490 & 8298490 & 0 & 100.0000 & 20 & 0.0000 & nan & nan & nan & nan & nan & nan & nan & nan & nan & nan & nan & nan & nan & nan & nan & nan & nan & nan & nan & nan & nan & nan & nan & \{'MG': 8297908, 'SP': 236, 'GO': 73, 'BA': 61, 'RJ': 44\} & 8297908.00 & 2.0000 & 2.0000 & 2.0000 & 0.0000 & 0.0000 & MG & 8297908.00 \\
qtd & int64 & NUMÉRICO & 8298490 & 8298490 & 0 & 100.0000 & 1410 & 0.0200 & 1.0000 & 10676.00 & 10675.00 & [1, 10676] & 18.8849 & 11.0000 & 1.0000 & 32.4566 & 1053.43 & 171.8653 & 29.1219 & 3940.81 & 4.0000 & 11.0000 & 24.0000 & 20.0000 & -26.0000 & 54.0000 & [4.0, 24.0] & 156716336.00 & nan & 1.0000 & 60.0000 & nan & nan & nan & nan & nan & nan & nan & nan & nan \\
\hline
\end{tabular}
\begin{flushleft}
\small Fonte: \textit{informacoes\_dados/bd\_gta\_dentro\_mg202505091607\_analise.csv}
\end{flushleft}
\end{table}


% \subsection{Análise por Faixa Etária}
% \begin{table}[H]
\centering
\small
\caption{Bd Gta Dentro Mg Faixa Etaria202505091602 Analise}
\label{tab:11_informacoes_dados}
\begin{tabular}{|l|r|r|r|r|r|r|r|r|r|r|r|r|r|r|r|r|r|r|r|r|r|r|r|r|r|r|r|r|r|r|r|r|r|r|r|r|r|r|r|r|}
\hline
\textbf{coluna} & \textbf{tipo_python} & \textbf{categoria_tipo} & \textbf{valores_totais} & \textbf{valores_validos} & \textbf{valores_nulos} & \textbf{percentual_preenchimento} & \textbf{valores_unicos} & \textbf{taxa_unica} & \textbf{minimo} & \textbf{maximo} & \textbf{amplitude} & \textbf{intervalo_textual} & \textbf{media} & \textbf{mediana} & \textbf{moda} & \textbf{desvio_padrao} & \textbf{variancia} & \textbf{coeficiente_variacao_pct} & \textbf{assimetria} & \textbf{curtose} & \textbf{quartil_25} & \textbf{quartil_50} & \textbf{quartil_75} & \textbf{amplitude_interquartil} & \textbf{limite_inferior_iqr} & \textbf{limite_superior_iqr} & \textbf{range_iqr} & \textbf{soma} & \textbf{produto} & \textbf{percentil_5} & \textbf{percentil_95} & \textbf{valores_top_5} & \textbf{frequencia_top} & \textbf{comprimento_min} & \textbf{comprimento_max} & \textbf{comprimento_medio} & \textbf{comprimento_desvio} & \textbf{diversidade} & \textbf{valor_mais_frequente} & \textbf{frequencia_maxima} \\
\hline
id_gta & int64 & NUMÉRICO & 11838896 & 11838896 & 0 & 100.0000 & 8301177 & 70.1200 & 6.0000 & 18222380.00 & 18222374.00 & [6, 18222380] & 8702163.76 & 8591654.50 & 1961.00 & 5384438.29 & 28992175687617.66 & 61.8747 & 0.0838 & -1.2485 & 3750167.50 & 8591654.50 & 13356966.25 & 9606798.75 & -10660030.62 & 27767164.38 & [3750167.5, 13356966.25] & 103024011776938.00 & nan & 765016.00 & 17280631.00 & nan & nan & nan & nan & nan & nan & nan & nan & nan \\
dt_emissao_gta & object & CATEGÓRICO & 11838896 & 11838896 & 0 & 100.0000 & 8267864 & 69.8400 & nan & nan & nan & nan & nan & nan & nan & nan & nan & nan & nan & nan & nan & nan & nan & nan & nan & nan & nan & nan & nan & nan & nan & {'2012-04-19 00:00:00.0': 2777, '2012-04-18 00:00:00.0': 2041, '2012-08-30 00:00:00.0': 619, '2019-06-01 00:00:00.0': 519, '2013-09-16 00:00:00.0': 470} & 2777.00 & 21.0000 & 23.0000 & 22.8800 & 0.3600 & 0.6984 & 2012-04-19 00:00:00.0 & 2777.00 \\
ds_faixa_etaria & object & CATEGÓRICO & 11838896 & 11838896 & 0 & 100.0000 & 4 & 0.0000 & nan & nan & nan & nan & nan & nan & nan & nan & nan & nan & nan & nan & nan & nan & nan & nan & nan & nan & nan & nan & nan & nan & nan & {'De 0 até 12 meses': 3884463, 'Acima de 36 meses': 3007839, 'De 13 até 24 meses': 2944075, 'De 25 até 36 meses': 2002519} & 3884463.00 & 17.0000 & 18.0000 & 17.4200 & 0.4900 & 0.0000 & De 0 até 12 meses & 3884463.00 \\
tp_sexo & object & CATEGÓRICO & 11838896 & 11838896 & 0 & 100.0000 & 2 & 0.0000 & nan & nan & nan & nan & nan & nan & nan & nan & nan & nan & nan & nan & nan & nan & nan & nan & nan & nan & nan & nan & nan & nan & nan & {'F': 6059997, 'M': 5778899} & 6059997.00 & 1.0000 & 1.0000 & 1.0000 & 0.0000 & 0.0000 & F & 6059997.00 \\
qt_animais & int64 & NUMÉRICO & 11838896 & 11838896 & 0 & 100.0000 & 1135 & 0.0100 & 1.0000 & 7282.00 & 7281.00 & [1, 7282] & 13.2381 & 6.0000 & 1.0000 & 23.2619 & 541.1140 & 175.7186 & 24.2133 & 3187.91 & 2.0000 & 6.0000 & 17.0000 & 15.0000 & -20.5000 & 39.5000 & [2.0, 17.0] & 156724866.00 & nan & 1.0000 & 45.0000 & nan & nan & nan & nan & nan & nan & nan & nan & nan \\
\hline
\end{tabular}
\begin{flushleft}
\small Fonte: \textit{informacoes_dados/bd_gta_dentro_mg_faixa_etaria202505091602_analise.csv}
\end{flushleft}
\end{table}


% \subsection{Análise OE}
% \begin{table}[H]
\centering
\small
\caption{Bd Gta Oe202505091618 Analise}
\label{tab:12_informacoes_dados}
\begin{tabular}{|l|r|r|r|r|r|r|r|r|r|r|r|r|r|r|r|r|r|r|r|r|r|r|r|r|r|r|r|r|r|r|r|r|r|r|r|r|r|r|r|r|}
\hline
\textbf{coluna} & \textbf{tipo\_python} & \textbf{categoria\_tipo} & \textbf{valores\_totais} & \textbf{valores\_validos} & \textbf{valores\_nulos} & \textbf{percentual\_preenchimento} & \textbf{valores\_unicos} & \textbf{taxa\_unica} & \textbf{minimo} & \textbf{maximo} & \textbf{amplitude} & \textbf{intervalo\_textual} & \textbf{media} & \textbf{mediana} & \textbf{moda} & \textbf{desvio\_padrao} & \textbf{variancia} & \textbf{coeficiente\_variacao\_pct} & \textbf{assimetria} & \textbf{curtose} & \textbf{quartil\_25} & \textbf{quartil\_50} & \textbf{quartil\_75} & \textbf{amplitude\_interquartil} & \textbf{limite\_inferior\_iqr} & \textbf{limite\_superior\_iqr} & \textbf{range\_iqr} & \textbf{soma} & \textbf{produto} & \textbf{percentil\_5} & \textbf{percentil\_95} & \textbf{valores\_top\_5} & \textbf{frequencia\_top} & \textbf{comprimento\_min} & \textbf{comprimento\_max} & \textbf{comprimento\_medio} & \textbf{comprimento\_desvio} & \textbf{diversidade} & \textbf{valor\_mais\_frequente} & \textbf{frequencia\_maxima} \\
\hline
nr\_gta & int64 & NUMÉRICO & 64677 & 64677 & 0 & 100.0000 & 62643 & 96.8600 & 23.0000 & 9831176.00 & 9831153.00 & [23, 9831176] & 482380.80 & 469824.00 & 281508.00 & 310018.25 & 96111313333.03 & 64.2684 & 2.2057 & 49.4139 & 222068.00 & 469824.00 & 744596.00 & 522528.00 & -561724.00 & 1528388.00 & [222068.0, 744596.0] & 31198943273.00 & nan & 40916.40 & 946428.00 & nan & nan & nan & nan & nan & nan & nan & nan & nan \\
nr\_serie & object & CATEGÓRICO & 64677 & 64677 & 0 & 100.0000 & 31 & 0.0500 & nan & nan & nan & nan & nan & nan & nan & nan & nan & nan & nan & nan & nan & nan & nan & nan & nan & nan & nan & nan & nan & nan & nan & \{'D ': 5778, 'N ': 4817, 'O ': 4437, 'K ': 4209, 'P ': 4032\} & 5778.00 & 2.0000 & 2.0000 & 2.0000 & 0.0000 & 0.0005 & D  & 5778.00 \\
id\_gta\_origem\_outro\_estado & int64 & NUMÉRICO & 64677 & 64677 & 0 & 100.0000 & 64677 & 100.0000 & 9.0000 & 358416.00 & 358407.00 & [9, 358416] & 159902.06 & 133004.00 & 9.0000 & 109783.03 & 12052314649.32 & 68.6564 & 0.2576 & -1.2628 & 60873.00 & 133004.00 & 259844.00 & 198971.00 & -237583.50 & 558300.50 & [60873.0, 259844.0] & 10341985763.00 & nan & 9156.40 & 339072.00 & nan & nan & nan & nan & nan & nan & nan & nan & nan \\
dt\_emissao\_gta & object & CATEGÓRICO & 64677 & 64677 & 0 & 100.0000 & 4108 & 6.3500 & nan & nan & nan & nan & nan & nan & nan & nan & nan & nan & nan & nan & nan & nan & nan & nan & nan & nan & nan & nan & nan & nan & nan & \{'2025-01-14 00:00:00.0': 86, '2023-09-01 00:00:00.0': 84, '2023-12-15 00:00:00.0': 79, '2023-02-13 00:00:00.0': 77, '2022-10-21 00:00:00.0': 77\} & 86.0000 & 21.0000 & 21.0000 & 21.0000 & 0.0000 & 0.0635 & 2025-01-14 00:00:00.0 & 86.0000 \\
cod\_produtor\_origem & int64 & NUMÉRICO & 64677 & 64677 & 0 & 100.0000 & 18472 & 28.5600 & 13.0000 & 80229.00 & 80216.00 & [13, 80229] & 26416.03 & 16921.00 & 72132.00 & 25882.43 & 669900415.65 & 97.9800 & 0.8928 & -0.6571 & 5487.00 & 16921.00 & 42864.00 & 37377.00 & -50578.50 & 98929.50 & [5487.0, 42864.0] & 1708509566.00 & nan & 442.0000 & 76423.00 & nan & nan & nan & nan & nan & nan & nan & nan & nan \\
cod\_produtor\_destino & int64 & NUMÉRICO & 64677 & 64677 & 0 & 100.0000 & 8918 & 13.7900 & 427.0000 & 3392697.00 & 3392270.00 & [427, 3392697] & 1237879.27 & 395460.00 & 23548.00 & 1419115.02 & 2013887435586.21 & 114.6408 & 0.6554 & -1.5242 & 107458.00 & 395460.00 & 3135779.00 & 3028321.00 & -4435023.50 & 7678260.50 & [107458.0, 3135779.0] & 80062317525.00 & nan & 18578.60 & 3312093.00 & nan & nan & nan & nan & nan & nan & nan & nan & nan \\
id\_tipo\_finalidade\_gta & int64 & NUMÉRICO & 64677 & 64677 & 0 & 100.0000 & 12 & 0.0200 & 2.0000 & 37.0000 & 35.0000 & [2, 37] & 4.6793 & 2.0000 & 2.0000 & 6.6169 & 43.7827 & 141.4079 & 3.3955 & 11.8045 & 2.0000 & 2.0000 & 3.0000 & 1.0000 & 0.5000 & 4.5000 & [2.0, 3.0] & 302641.00 & nan & 2.0000 & 18.0000 & nan & nan & nan & nan & nan & nan & nan & nan & nan \\
municipio\_origem & object & CATEGÓRICO & 64677 & 64677 & 0 & 100.0000 & 2121 & 3.2800 & nan & nan & nan & nan & nan & nan & nan & nan & nan & nan & nan & nan & nan & nan & nan & nan & nan & nan & nan & nan & nan & nan & nan & \{'SANTA CRUZ DO XINGU': 2405, 'JUSSARA': 991, 'GURUPI': 978, 'PARANAÍBA': 795, 'ITARUMÃ': 628\} & 2405.00 & 3.0000 & 32.0000 & 11.1800 & 5.1500 & 0.0328 & SANTA CRUZ DO XINGU & 2405.00 \\
cod\_municipio\_origem & int64 & NUMÉRICO & 64677 & 64677 & 0 & 100.0000 & 2156 & 3.3300 & 1100015.00 & 5300108.00 & 4200093.00 & [1100015, 5300108] & 3815396.90 & 3538907.00 & 5107743.00 & 1213807.84 & 1473329461760.28 & 31.8134 & -0.4116 & -0.8023 & 3200409.00 & 3538907.00 & 5106315.00 & 1905906.00 & 341550.00 & 7965174.00 & [3200409.0, 5106315.0] & 246768425498.00 & nan & 1506138.00 & 5214051.00 & nan & nan & nan & nan & nan & nan & nan & nan & nan \\
estado\_origem & object & CATEGÓRICO & 64677 & 64677 & 0 & 100.0000 & 24 & 0.0400 & nan & nan & nan & nan & nan & nan & nan & nan & nan & nan & nan & nan & nan & nan & nan & nan & nan & nan & nan & nan & nan & nan & nan & \{'SP': 14320, 'GO': 11257, 'MT': 8098, 'BA': 6859, 'MS': 4482\} & 14320.00 & 2.0000 & 2.0000 & 2.0000 & 0.0000 & 0.0004 & SP & 14320.00 \\
municipio\_destino & object & CATEGÓRICO & 64677 & 64677 & 0 & 100.0000 & 775 & 1.2000 & nan & nan & nan & nan & nan & nan & nan & nan & nan & nan & nan & nan & nan & nan & nan & nan & nan & nan & nan & nan & nan & nan & nan & \{'UBERABA': 6122, 'FRUTAL': 3109, 'CAPINÓPOLIS': 2920, 'CARNEIRINHO': 2249, 'GUARDA-MOR': 1917\} & 6122.00 & 3.0000 & 30.0000 & 10.7200 & 4.6500 & 0.0120 & UBERABA & 6122.00 \\
cod\_municipio\_destino & int64 & NUMÉRICO & 64677 & 64677 & 0 & 100.0000 & 775 & 1.2000 & 3100104.00 & 3172202.00 & 72098.00 & [3100104, 3172202] & 3138344.58 & 3137502.00 & 3170107.00 & 22250.40 & 495080304.66 & 0.7090 & 0.0582 & -1.3174 & 3115805.00 & 3137502.00 & 3159803.00 & 43998.00 & 3049808.00 & 3225800.00 & [3115805.0, 3159803.0] & 202978712602.00 & nan & 3104700.00 & 3170206.00 & nan & nan & nan & nan & nan & nan & nan & nan & nan \\
estado\_destino & object & CATEGÓRICO & 64677 & 64677 & 0 & 100.0000 & 1 & 0.0000 & nan & nan & nan & nan & nan & nan & nan & nan & nan & nan & nan & nan & nan & nan & nan & nan & nan & nan & nan & nan & nan & nan & nan & \{'MG': 64677\} & 64677.00 & 2.0000 & 2.0000 & 2.0000 & 0.0000 & 0.0000 & MG & 64677.00 \\
qtd & int64 & NUMÉRICO & 64677 & 64677 & 0 & 100.0000 & 236 & 0.3600 & 1.0000 & 1000.00 & 999.0000 & [1, 1000] & 37.1705 & 25.0000 & 1.0000 & 41.2141 & 1698.60 & 110.8785 & 4.2369 & 45.7795 & 5.0000 & 25.0000 & 63.0000 & 58.0000 & -82.0000 & 150.0000 & [5.0, 63.0] & 2404077.00 & nan & 1.0000 & 100.0000 & nan & nan & nan & nan & nan & nan & nan & nan & nan \\
\hline
\end{tabular}
\begin{flushleft}
\small Fonte: \textit{informacoes\_dados/bd\_gta\_oe202505091618\_analise.csv}
\end{flushleft}
\end{table}


% \subsection{Análise OE por Faixa Etária}
% \begin{table}[H]
\centering
\small
\caption{Bd Gta Oe Faixa Etaria202505091630 Analise}
\label{tab:13_informacoes_dados}
\begin{tabular}{|l|r|r|r|r|r|r|r|r|r|r|r|r|r|r|r|r|r|r|r|r|r|r|r|r|r|r|r|r|r|r|r|r|r|r|r|r|r|r|r|r|}
\hline
\textbf{coluna} & \textbf{tipo_python} & \textbf{categoria_tipo} & \textbf{valores_totais} & \textbf{valores_validos} & \textbf{valores_nulos} & \textbf{percentual_preenchimento} & \textbf{valores_unicos} & \textbf{taxa_unica} & \textbf{minimo} & \textbf{maximo} & \textbf{amplitude} & \textbf{intervalo_textual} & \textbf{media} & \textbf{mediana} & \textbf{moda} & \textbf{desvio_padrao} & \textbf{variancia} & \textbf{coeficiente_variacao_pct} & \textbf{assimetria} & \textbf{curtose} & \textbf{quartil_25} & \textbf{quartil_50} & \textbf{quartil_75} & \textbf{amplitude_interquartil} & \textbf{limite_inferior_iqr} & \textbf{limite_superior_iqr} & \textbf{range_iqr} & \textbf{soma} & \textbf{produto} & \textbf{percentil_5} & \textbf{percentil_95} & \textbf{valores_top_5} & \textbf{frequencia_top} & \textbf{comprimento_min} & \textbf{comprimento_max} & \textbf{comprimento_medio} & \textbf{comprimento_desvio} & \textbf{diversidade} & \textbf{valor_mais_frequente} & \textbf{frequencia_maxima} \\
\hline
id_gta_origem_outro_estado & int64 & NUMÉRICO & 87560 & 87560 & 0 & 100.0000 & 64678 & 73.8700 & 9.0000 & 358416.00 & 358407.00 & [9, 358416] & 155633.34 & 127725.50 & 30419.00 & 109635.10 & 12019856152.32 & 70.4445 & 0.3082 & -1.2340 & 56372.00 & 127725.50 & 254649.25 & 198277.25 & -241043.88 & 552065.12 & [56372.0, 254649.25] & 13627255607.00 & nan & 8185.00 & 337700.20 & nan & nan & nan & nan & nan & nan & nan & nan & nan \\
dt_emissao_gta & object & CATEGÓRICO & 87560 & 87560 & 0 & 100.0000 & 4108 & 4.6900 & nan & nan & nan & nan & nan & nan & nan & nan & nan & nan & nan & nan & nan & nan & nan & nan & nan & nan & nan & nan & nan & nan & nan & {'2022-10-21 00:00:00.0': 122, '2024-09-18 00:00:00.0': 106, '2024-06-14 00:00:00.0': 106, '2025-01-14 00:00:00.0': 100, '2022-11-22 00:00:00.0': 98} & 122.0000 & 21.0000 & 21.0000 & 21.0000 & 0.0000 & 0.0469 & 2022-10-21 00:00:00.0 & 122.0000 \\
ds_faixa_etaria & object & CATEGÓRICO & 87560 & 87560 & 0 & 100.0000 & 4 & 0.0000 & nan & nan & nan & nan & nan & nan & nan & nan & nan & nan & nan & nan & nan & nan & nan & nan & nan & nan & nan & nan & nan & nan & nan & {'De 13 até 24 meses': 28767, 'De 0 até 12 meses': 25939, 'Acima de 36 meses': 16527, 'De 25 até 36 meses': 16327} & 28767.00 & 17.0000 & 18.0000 & 17.5200 & 0.5000 & 0.0000 & De 13 até 24 meses & 28767.00 \\
tp_sexo & object & CATEGÓRICO & 87560 & 87560 & 0 & 100.0000 & 2 & 0.0000 & nan & nan & nan & nan & nan & nan & nan & nan & nan & nan & nan & nan & nan & nan & nan & nan & nan & nan & nan & nan & nan & nan & nan & {'M': 49689, 'F': 37871} & 49689.00 & 1.0000 & 1.0000 & 1.0000 & 0.0000 & 0.0000 & M & 49689.00 \\
qt_animais & int64 & NUMÉRICO & 87560 & 87560 & 0 & 100.0000 & 219 & 0.2500 & 1.0000 & 1000.00 & 999.0000 & [1, 1000] & 27.4564 & 13.0000 & 1.0000 & 35.2175 & 1240.27 & 128.2670 & 3.4830 & 33.9414 & 2.0000 & 13.0000 & 47.0000 & 45.0000 & -65.5000 & 114.5000 & [2.0, 47.0] & 2404085.00 & nan & 1.0000 & 90.0000 & nan & nan & nan & nan & nan & nan & nan & nan & nan \\
\hline
\end{tabular}
\begin{flushleft}
\small Fonte: \textit{informacoes_dados/bd_gta_oe_faixa_etaria202505091630_analise.csv}
\end{flushleft}
\end{table}


% \subsection{Análise com Distâncias}
% \begin{table}[H]
\centering
\small
\caption{Gtas Com Distancias Analise}
\label{tab:14_informacoes_dados}
\begin{tabular}{|l|r|r|r|r|r|r|r|r|r|r|r|r|r|r|r|r|r|r|r|r|r|r|r|r|r|r|r|r|r|r|r|r|r|r|r|r|r|r|r|r|}
\hline
\textbf{coluna} & \textbf{tipo\_python} & \textbf{categoria\_tipo} & \textbf{valores\_totais} & \textbf{valores\_validos} & \textbf{valores\_nulos} & \textbf{percentual\_preenchimento} & \textbf{valores\_unicos} & \textbf{taxa\_unica} & \textbf{minimo} & \textbf{maximo} & \textbf{amplitude} & \textbf{intervalo\_textual} & \textbf{media} & \textbf{mediana} & \textbf{moda} & \textbf{desvio\_padrao} & \textbf{variancia} & \textbf{coeficiente\_variacao\_pct} & \textbf{assimetria} & \textbf{curtose} & \textbf{quartil\_25} & \textbf{quartil\_50} & \textbf{quartil\_75} & \textbf{amplitude\_interquartil} & \textbf{limite\_inferior\_iqr} & \textbf{limite\_superior\_iqr} & \textbf{range\_iqr} & \textbf{soma} & \textbf{produto} & \textbf{percentil\_5} & \textbf{percentil\_95} & \textbf{valores\_top\_5} & \textbf{frequencia\_top} & \textbf{comprimento\_min} & \textbf{comprimento\_max} & \textbf{comprimento\_medio} & \textbf{comprimento\_desvio} & \textbf{diversidade} & \textbf{valor\_mais\_frequente} & \textbf{frequencia\_maxima} \\
\hline
nr\_gta & int64 & NUMÉRICO & 8363167 & 8363167 & 0 & 100.0000 & 1000067 & 11.9600 & 1.0000 & 66760816.00 & 66760815.00 & [1, 66760816] & 495055.46 & 493920.00 & 1.0000 & 293024.49 & 85863349061.85 & 59.1902 & 3.7872 & 756.2359 & 241777.00 & 493920.00 & 746996.00 & 505219.00 & -516051.50 & 1504824.50 & [241777.0, 746996.0] & 4140231475884.00 & nan & 47705.30 & 949077.00 & nan & nan & nan & nan & nan & nan & nan & nan & nan \\
nr\_serie & object & CATEGÓRICO & 8363167 & 8363167 & 0 & 100.0000 & 51 & 0.0000 & nan & nan & nan & nan & nan & nan & nan & nan & nan & nan & nan & nan & nan & nan & nan & nan & nan & nan & nan & nan & nan & nan & nan & \{'B': 538325, 'C': 532592, 'A': 526540, 'L': 488731, 'D': 486189\} & 538325.00 & 1.0000 & 2.0000 & 1.0100 & 0.0900 & 0.0000 & B & 538325.00 \\
id\_gta & int64 & NUMÉRICO & 8363167 & 8363167 & 0 & 100.0000 & 8330547 & 99.6100 & 6.0000 & 18222380.00 & 18222374.00 & [6, 18222380] & 8816874.99 & 8791998.00 & 10.0000 & 5417607.99 & 29350476318791.89 & 61.4459 & 0.0456 & -1.2544 & 3871568.00 & 8791998.00 & 13542842.00 & 9671274.00 & -10635343.00 & 28049753.00 & [3871568.0, 13542842.0] & 73736997991033.00 & nan & 691953.60 & 17317988.70 & nan & nan & nan & nan & nan & nan & nan & nan & nan \\
dt\_emissao\_gta & object & CATEGÓRICO & 8363167 & 8363167 & 0 & 100.0000 & 8268228 & 98.8600 & nan & nan & nan & nan & nan & nan & nan & nan & nan & nan & nan & nan & nan & nan & nan & nan & nan & nan & nan & nan & nan & nan & nan & \{'2012-04-19 00:00:00.0': 1753, '2012-04-18 00:00:00.0': 1297, '2012-08-30 00:00:00.0': 456, '2019-06-01 00:00:00.0': 396, '2013-09-16 00:00:00.0': 338\} & 1753.00 & 21.0000 & 23.0000 & 22.8700 & 0.4000 & 0.9886 & 2012-04-19 00:00:00.0 & 1753.00 \\
cod\_produtor\_origem & int64 & NUMÉRICO & 8363167 & 8363167 & 0 & 100.0000 & 464877 & 5.5600 & 1.0000 & 3393952.00 & 3393951.00 & [1, 3393952] & 882199.67 & 290281.00 & 77922.00 & 1237197.61 & 1530657926143.35 & 140.2401 & 1.2906 & -0.2632 & 110921.00 & 290281.00 & 491227.00 & 380306.00 & -459538.00 & 1061686.00 & [110921.0, 491227.0] & 7377983181752.00 & nan & 19506.00 & 3252072.00 & nan & nan & nan & nan & nan & nan & nan & nan & nan \\
cod\_produtor\_destino & int64 & NUMÉRICO & 8363167 & 8363167 & 0 & 100.0000 & 416822 & 4.9800 & 2.0000 & 3394030.00 & 3394028.00 & [2, 3394030] & 1041702.11 & 322326.00 & 17348.00 & 1331310.52 & 1772387712296.53 & 127.8015 & 0.9697 & -1.0051 & 112684.00 & 322326.00 & 3071041.00 & 2958357.00 & -4324851.50 & 7508576.50 & [112684.0, 3071041.0] & 8711928696548.00 & nan & 18601.00 & 3274939.00 & nan & nan & nan & nan & nan & nan & nan & nan & nan \\
id\_tipo\_finalidade\_gta & int64 & NUMÉRICO & 8363167 & 8363167 & 0 & 100.0000 & 14 & 0.0000 & 2.0000 & 37.0000 & 35.0000 & [2, 37] & 13.6278 & 8.0000 & 2.0000 & 11.7889 & 138.9783 & 86.5061 & 0.4822 & -1.2964 & 2.0000 & 8.0000 & 30.0000 & 28.0000 & -40.0000 & 72.0000 & [2.0, 30.0] & 113971845.00 & nan & 2.0000 & 30.0000 & nan & nan & nan & nan & nan & nan & nan & nan & nan \\
ds\_meio\_transporte & object & CATEGÓRICO & 8363167 & 8298490 & 64677 & 99.2300 & 5 & 0.0000 & nan & nan & nan & nan & nan & nan & nan & nan & nan & nan & nan & nan & nan & nan & nan & nan & nan & nan & nan & nan & nan & nan & nan & \{'RODOVIÁRIO': 7322115, 'A PÉ': 968445, 'FERROVIÁRIO': 4297, 'MARÍTIMO/FLUVIAL': 2938, 'AÉREO': 695\} & 7322115.00 & 4.0000 & 16.0000 & 9.3000 & 1.9300 & 0.0000 & RODOVIÁRIO & 7322115.00 \\
cod\_municipio\_origem & int64 & NUMÉRICO & 8363167 & 8363167 & 0 & 100.0000 & 3009 & 0.0400 & 1100015.00 & 5300108.00 & 4200093.00 & [1100015, 5300108] & 3142885.16 & 3137809.00 & 3127107.00 & 123939.20 & 15360924183.64 & 3.9435 & 10.3163 & 226.1990 & 3120805.00 & 3137809.00 & 3154309.00 & 33504.00 & 3070549.00 & 3204565.00 & [3120805.0, 3154309.0] & 26284473430623.00 & nan & 3103504.00 & 3170206.00 & nan & nan & nan & nan & nan & nan & nan & nan & nan \\
nome\_municipio\_origem & object & CATEGÓRICO & 8363167 & 8363167 & 0 & 100.0000 & 2955 & 0.0400 & nan & nan & nan & nan & nan & nan & nan & nan & nan & nan & nan & nan & nan & nan & nan & nan & nan & nan & nan & nan & nan & nan & nan & \{'FRUTAL': 284929, 'UBERLÂNDIA': 236250, 'PATOS DE MINAS': 198900, 'PRATA': 166060, 'UBERABA': 163394\} & 284929.00 & 3.0000 & 32.0000 & 10.7600 & 4.7300 & 0.0004 & FRUTAL & 284929.00 \\
estado\_origem & object & CATEGÓRICO & 8363167 & 8363167 & 0 & 100.0000 & 25 & 0.0000 & nan & nan & nan & nan & nan & nan & nan & nan & nan & nan & nan & nan & nan & nan & nan & nan & nan & nan & nan & nan & nan & nan & nan & \{'MG': 8298490, 'SP': 14320, 'GO': 11257, 'MT': 8098, 'BA': 6859\} & 8298490.00 & 2.0000 & 2.0000 & 2.0000 & 0.0000 & 0.0000 & MG & 8298490.00 \\
cod\_municipio\_destino & int64 & NUMÉRICO & 8363167 & 8363167 & 0 & 100.0000 & 1138 & 0.0100 & 1100122.00 & 5300108.00 & 4199986.00 & [1100122, 5300108] & 3137352.12 & 3137106.00 & 3127107.00 & 22949.75 & 526691229.34 & 0.7315 & 10.2556 & 1140.73 & 3119302.00 & 3137106.00 & 3153400.00 & 34098.00 & 3068155.00 & 3204547.00 & [3119302.0, 3153400.0] & 26238199692266.00 & nan & 3103504.00 & 3170206.00 & nan & nan & nan & nan & nan & nan & nan & nan & nan \\
nome\_municipio\_destino & object & CATEGÓRICO & 8363167 & 8363167 & 0 & 100.0000 & 1136 & 0.0100 & nan & nan & nan & nan & nan & nan & nan & nan & nan & nan & nan & nan & nan & nan & nan & nan & nan & nan & nan & nan & nan & nan & nan & \{'FRUTAL': 210168, 'PRATA': 207093, 'CAMPINA VERDE': 192091, 'UBERABA': 140418, 'UNAÍ': 133672\} & 210168.00 & 3.0000 & 30.0000 & 10.9500 & 4.8000 & 0.0001 & FRUTAL & 210168.00 \\
estado\_destino & object & CATEGÓRICO & 8363167 & 8363167 & 0 & 100.0000 & 20 & 0.0000 & nan & nan & nan & nan & nan & nan & nan & nan & nan & nan & nan & nan & nan & nan & nan & nan & nan & nan & nan & nan & nan & nan & nan & \{'MG': 8362585, 'SP': 236, 'GO': 73, 'BA': 61, 'RJ': 44\} & 8362585.00 & 2.0000 & 2.0000 & 2.0000 & 0.0000 & 0.0000 & MG & 8362585.00 \\
qtd & int64 & NUMÉRICO & 8363167 & 8363167 & 0 & 100.0000 & 1410 & 0.0200 & 1.0000 & 10676.00 & 10675.00 & [1, 10676] & 19.0263 & 11.0000 & 1.0000 & 32.5728 & 1060.99 & 171.1986 & 28.6642 & 3855.42 & 4.0000 & 11.0000 & 24.0000 & 20.0000 & -26.0000 & 54.0000 & [4.0, 24.0] & 159120413.00 & nan & 1.0000 & 60.0000 & nan & nan & nan & nan & nan & nan & nan & nan & nan \\
Acima de 36 meses & int64 & NUMÉRICO & 8363167 & 8363167 & 0 & 100.0000 & 711 & 0.0100 & 0.0000 & 5861.00 & 5861.00 & [0, 5861] & 3.8626 & 0.0000 & 0.0000 & 13.3565 & 178.3957 & 345.7921 & 38.9686 & 7956.51 & 0.0000 & 0.0000 & 2.0000 & 2.0000 & -3.0000 & 5.0000 & [0.0, 2.0] & 32303374.00 & nan & 0.0000 & 20.0000 & nan & nan & nan & nan & nan & nan & nan & nan & nan \\
De 0 até 12 meses & int64 & NUMÉRICO & 8363167 & 8363167 & 0 & 100.0000 & 756 & 0.0100 & 0.0000 & 7282.00 & 7282.00 & [0, 7282] & 6.3298 & 0.0000 & 0.0000 & 17.6055 & 309.9547 & 278.1353 & 36.0187 & 9045.25 & 0.0000 & 0.0000 & 5.0000 & 5.0000 & -7.5000 & 12.5000 & [0.0, 5.0] & 52937548.00 & nan & 0.0000 & 33.0000 & nan & nan & nan & nan & nan & nan & nan & nan & nan \\
De 13 até 24 meses & int64 & NUMÉRICO & 8363167 & 8363167 & 0 & 100.0000 & 858 & 0.0100 & 0.0000 & 3000.00 & 3000.00 & [0, 3000] & 5.5517 & 0.0000 & 0.0000 & 17.5086 & 306.5510 & 315.3721 & 19.9571 & 1319.13 & 0.0000 & 0.0000 & 3.0000 & 3.0000 & -4.5000 & 7.5000 & [0.0, 3.0] & 46430010.00 & nan & 0.0000 & 30.0000 & nan & nan & nan & nan & nan & nan & nan & nan & nan \\
De 25 até 36 meses & int64 & NUMÉRICO & 8363167 & 8363167 & 0 & 100.0000 & 752 & 0.0100 & 0.0000 & 3968.00 & 3968.00 & [0, 3968] & 3.2822 & 0.0000 & 0.0000 & 13.6300 & 185.7770 & 415.2719 & 32.6136 & 3991.20 & 0.0000 & 0.0000 & 0.0000 & 0.0000 & 0.0000 & 0.0000 & [0.0, 0.0] & 27449481.00 & nan & 0.0000 & 20.0000 & nan & nan & nan & nan & nan & nan & nan & nan & nan \\
distancia\_km & float64 & NUMÉRICO & 8363167 & 8295186 & 67981 & 99.1900 & 144885 & 1.7500 & 0.0000 & 5321.43 & 5321.43 & [0.0, 5321.4315] & 87.8809 & 29.7449 & 0.0000 & 167.1733 & 27946.90 & 190.2272 & 3.5667 & 17.4131 & 0.0000 & 29.7449 & 91.6114 & 91.6114 & -137.4171 & 229.0285 & [0.0, 91.6114] & 728988062.28 & nan & 0.0000 & 426.2785 & nan & nan & nan & nan & nan & nan & nan & nan & nan \\
duracao\_min & float64 & NUMÉRICO & 8363167 & 8295186 & 67981 & 99.1900 & 125553 & 1.5100 & 0.0000 & 19076.71 & 19076.71 & [0.0, 19076.71] & 100.2290 & 33.9633 & 0.0000 & 285.8673 & 81720.11 & 285.2142 & 21.3238 & 774.0778 & 0.0000 & 33.9633 & 99.9017 & 99.9017 & -149.8525 & 249.7542 & [0.0, 99.9017] & 831418087.14 & nan & 0.0000 & 449.9383 & nan & nan & nan & nan & nan & nan & nan & nan & nan \\
\hline
\end{tabular}
\begin{flushleft}
\small Fonte: \textit{informacoes\_dados/gtas\_com\_distancias\_analise.csv}
\end{flushleft}
\end{table}


% \subsection{Matriz de Distâncias}
% \begin{table}[H]
\centering
\small
\caption{Matriz Distancias Analise}
\label{tab:15_informacoes_dados}
\begin{tabular}{|l|r|r|r|r|r|r|r|r|r|r|r|r|r|r|r|r|r|r|r|r|r|r|r|r|r|r|r|r|r|r|r|r|r|r|r|r|r|r|r|r|}
\hline
\textbf{coluna} & \textbf{tipo_python} & \textbf{categoria_tipo} & \textbf{valores_totais} & \textbf{valores_validos} & \textbf{valores_nulos} & \textbf{percentual_preenchimento} & \textbf{valores_unicos} & \textbf{taxa_unica} & \textbf{valores_top_5} & \textbf{frequencia_top} & \textbf{comprimento_min} & \textbf{comprimento_max} & \textbf{comprimento_medio} & \textbf{comprimento_desvio} & \textbf{diversidade} & \textbf{valor_mais_frequente} & \textbf{frequencia_maxima} & \textbf{minimo} & \textbf{maximo} & \textbf{amplitude} & \textbf{intervalo_textual} & \textbf{media} & \textbf{mediana} & \textbf{moda} & \textbf{desvio_padrao} & \textbf{variancia} & \textbf{coeficiente_variacao_pct} & \textbf{assimetria} & \textbf{curtose} & \textbf{quartil_25} & \textbf{quartil_50} & \textbf{quartil_75} & \textbf{amplitude_interquartil} & \textbf{limite_inferior_iqr} & \textbf{limite_superior_iqr} & \textbf{range_iqr} & \textbf{soma} & \textbf{produto} & \textbf{percentil_5} & \textbf{percentil_95} \\
\hline
municipio_origem & object & CATEGÓRICO & 4751210 & 4751210 & 0 & 100.0000 & 853 & 0.0200 & {'Abadia dos Dourados': 5570, 'Abaeté': 5570, 'Abre Campo': 5570, 'Acaiaca': 5570, 'Açucena': 5570} & 5570.00 & 3.0000 & 30.0000 & 11.9600 & 5.2100 & 0.0002 & Abadia dos Dourados & 5570.00 & nan & nan & nan & nan & nan & nan & nan & nan & nan & nan & nan & nan & nan & nan & nan & nan & nan & nan & nan & nan & nan & nan & nan \\
municipio_destino & object & CATEGÓRICO & 4751210 & 4751210 & 0 & 100.0000 & 5299 & 0.1100 & {'Bom Jesus': 4265, 'São Domingos': 4265, 'Santa Luzia': 3412, 'Santa Helena': 3412, 'Planalto': 3412} & 4265.00 & 3.0000 & 32.0000 & 11.6200 & 5.1800 & 0.0011 & Bom Jesus & 4265.00 & nan & nan & nan & nan & nan & nan & nan & nan & nan & nan & nan & nan & nan & nan & nan & nan & nan & nan & nan & nan & nan & nan & nan \\
dist_rod_km & float64 & NUMÉRICO & 4751210 & 4744386 & 6824 & 99.8600 & 4351712 & 91.7200 & nan & nan & nan & nan & nan & nan & nan & nan & nan & 0.0000 & 6083.52 & 6083.52 & [0.0, 6083.5158] & 1446.33 & 1385.71 & 0.0000 & 810.2920 & 656573.08 & 56.0238 & 0.8743 & 1.7009 & 825.4536 & 1385.71 & 1965.57 & 1140.12 & -884.7217 & 3675.75 & [825.4536, 1965.5704] & 6861968205.60 & nan & 310.1977 & 2707.31 \\
dist_rod_min & float64 & NUMÉRICO & 4751210 & 4744386 & 6824 & 99.8600 & 1466278 & 30.9100 & nan & nan & nan & nan & nan & nan & nan & nan & nan & 0.0000 & 23235.44 & 23235.44 & [0.0, 23235.4383] & 1285.82 & 1166.85 & 0.0000 & 1143.60 & 1307830.15 & 88.9394 & 9.0758 & 133.9071 & 698.7650 & 1166.85 & 1680.25 & 981.4817 & -773.4575 & 3152.47 & [698.765, 1680.2467] & 6100446677.91 & nan & 287.0267 & 2306.75 \\
dist_ped_km & float64 & NUMÉRICO & 4751210 & 4658376 & 92834 & 98.0500 & 3975828 & 85.3500 & nan & nan & nan & nan & nan & nan & nan & nan & nan & 0.0000 & 8121.13 & 8121.13 & [0.0, 8121.1269] & 1458.41 & 1405.29 & 0.0000 & 871.2587 & 759091.64 & 59.7402 & 1.7734 & 8.5815 & 825.8322 & 1405.29 & 1972.68 & 1146.85 & -894.4369 & 3692.95 & [825.8322, 1972.6783] & 6793839188.13 & nan & 299.4547 & 2616.96 \\
dist_ped_min & float64 & NUMÉRICO & 4751210 & 4658376 & 92834 & 98.0500 & 3837009 & 82.3700 & nan & nan & nan & nan & nan & nan & nan & nan & nan & 0.0000 & 89479.19 & 89479.19 & [0.0, 89479.19] & 17309.17 & 16741.41 & 0.0000 & 9988.41 & 99768314.31 & 57.7059 & 1.4370 & 6.3250 & 9892.93 & 16741.41 & 23482.23 & 13589.29 & -10491.01 & 43866.17 & [9892.9317, 23482.2262] & 80632634816.92 & nan & 3594.01 & 31343.62 \\
\hline
\end{tabular}
\begin{flushleft}
\small Fonte: \textit{informacoes_dados/matriz_distancias_analise.csv}
\end{flushleft}
\end{table}


% \subsection{Análise de Fraudes}
% \begin{table}[H]
\centering
\small
\caption{Possiveis Fraudes Analise}
\label{tab:16_informacoes_dados}
\begin{tabular}{|l|r|r|r|r|r|r|r|r|r|r|r|r|r|r|r|r|r|r|r|r|r|r|r|r|r|r|r|r|r|r|r|r|r|r|r|r|r|r|r|r|}
\hline
\textbf{coluna} & \textbf{tipo\_python} & \textbf{categoria\_tipo} & \textbf{valores\_totais} & \textbf{valores\_validos} & \textbf{valores\_nulos} & \textbf{percentual\_preenchimento} & \textbf{valores\_unicos} & \textbf{taxa\_unica} & \textbf{minimo} & \textbf{maximo} & \textbf{amplitude} & \textbf{intervalo\_textual} & \textbf{media} & \textbf{mediana} & \textbf{moda} & \textbf{desvio\_padrao} & \textbf{variancia} & \textbf{coeficiente\_variacao\_pct} & \textbf{assimetria} & \textbf{curtose} & \textbf{quartil\_25} & \textbf{quartil\_50} & \textbf{quartil\_75} & \textbf{amplitude\_interquartil} & \textbf{limite\_inferior\_iqr} & \textbf{limite\_superior\_iqr} & \textbf{range\_iqr} & \textbf{soma} & \textbf{produto} & \textbf{percentil\_5} & \textbf{percentil\_95} & \textbf{valores\_top\_5} & \textbf{frequencia\_top} & \textbf{comprimento\_min} & \textbf{comprimento\_max} & \textbf{comprimento\_medio} & \textbf{comprimento\_desvio} & \textbf{diversidade} & \textbf{valor\_mais\_frequente} & \textbf{frequencia\_maxima} \\
\hline
id\_gta & int64 & NUMÉRICO & 311 & 311 & 0 & 100.0000 & 311 & 100.0000 & 6216.00 & 18072282.00 & 18066066.00 & [6216, 18072282] & 8325406.80 & 8105213.00 & 6216.00 & 5461870.47 & 29832029049701.52 & 65.6048 & 0.1279 & -1.2378 & 3314671.50 & 8105213.00 & 13174684.50 & 9860013.00 & -11475348.00 & 27964704.00 & [3314671.5, 13174684.5] & 2589201514.00 & nan & 230558.50 & 17114282.50 & nan & nan & nan & nan & nan & nan & nan & nan & nan \\
dt\_emissao\_gta & object & CATEGÓRICO & 311 & 311 & 0 & 100.0000 & 311 & 100.0000 & nan & nan & nan & nan & nan & nan & nan & nan & nan & nan & nan & nan & nan & nan & nan & nan & nan & nan & nan & nan & nan & nan & nan & \{'2017-09-15 08:43:07.588': 1, '2021-06-16 08:12:09.641': 1, '2013-09-04 08:59:57.521': 1, '2016-12-07 09:29:34.062': 1, '2016-12-20 15:53:01.167': 1\} & 1.0000 & 23.0000 & 23.0000 & 23.0000 & 0.0000 & 1.0000 & 2017-09-15 08:43:07.588 & 1.0000 \\
cod\_produtor\_origem & int64 & NUMÉRICO & 311 & 311 & 0 & 100.0000 & 283 & 91.0000 & 736.0000 & 3359222.00 & 3358486.00 & [736, 3359222] & 980430.96 & 288691.00 & 344083.00 & 1328715.78 & 1765485626343.29 & 135.5236 & 1.0445 & -0.8576 & 73247.00 & 288691.00 & 3074223.50 & 3000976.50 & -4428217.75 & 7575688.25 & [73247.0, 3074223.5] & 304914028.00 & nan & 15650.50 & 3267098.00 & nan & nan & nan & nan & nan & nan & nan & nan & nan \\
cod\_produtor\_destino & int64 & NUMÉRICO & 311 & 311 & 0 & 100.0000 & 203 & 65.2700 & 2190.00 & 3310463.00 & 3308273.00 & [2190, 3310463] & 963562.58 & 290259.00 & 2190.00 & 1309625.96 & 1715120157637.03 & 135.9150 & 1.0816 & -0.7753 & 73024.00 & 290259.00 & 3051016.00 & 2977992.00 & -4393964.00 & 7518004.00 & [73024.0, 3051016.0] & 299667961.00 & nan & 16443.00 & 3259100.00 & nan & nan & nan & nan & nan & nan & nan & nan & nan \\
qtd & int64 & NUMÉRICO & 311 & 311 & 0 & 100.0000 & 210 & 67.5200 & 1.0000 & 4603.00 & 4602.00 & [1, 4603] & 350.0386 & 170.0000 & 1.0000 & 520.8562 & 271291.21 & 148.7997 & 3.4687 & 18.0603 & 40.0000 & 170.0000 & 421.5000 & 381.5000 & -532.2500 & 993.7500 & [40.0, 421.5] & 108862.00 & nan & 3.0000 & 1388.00 & nan & nan & nan & nan & nan & nan & nan & nan & nan \\
distancia\_km & float64 & NUMÉRICO & 311 & 311 & 0 & 100.0000 & 180 & 57.8800 & 0.0000 & 2469.95 & 2469.95 & [0.0, 2469.95] & 292.3097 & 55.0800 & 0.0000 & 449.9431 & 202448.77 & 153.9268 & 2.0502 & 4.7183 & 0.0000 & 55.0800 & 476.7350 & 476.7350 & -715.1025 & 1191.84 & [0.0, 476.735] & 90908.32 & nan & 0.0000 & 1165.69 & nan & nan & nan & nan & nan & nan & nan & nan & nan \\
data\_seconds & int64 & NUMÉRICO & 311 & 311 & 0 & 100.0000 & 311 & 100.0000 & 1336144677.00 & 1744465375.00 & 408320698.00 & [1336144677, 1744465375] & 1552026769.74 & 1559825115.00 & 1336144677.00 & 120163842.85 & 14439349127288274.00 & 7.7424 & -0.0844 & -1.2332 & 1443499943.00 & 1559825115.00 & 1660463184.00 & 216963241.00 & 1118055081.50 & 1985908045.50 & [1443499943.0, 1660463184.0] & 482680325388.00 & nan & 1357900595.00 & 1729721581.50 & nan & nan & nan & nan & nan & nan & nan & nan & nan \\
\hline
\end{tabular}
\begin{flushleft}
\small Fonte: \textit{informacoes\_dados/possiveis\_fraudes\_analise.csv}
\end{flushleft}
\end{table}


\end{document}

%% ============================================================================
%% NOTAS IMPORTANTES:
%% 
%% 1. Certifique-se que os caminhos dos arquivos \input{} estão corretos
%% 2. Execute: pdflatex exemplo_uso_tabelas.tex
%% 3. As tabelas são geradas automaticamente pelo notebook gerar_tabelas_latex.ipynb
%% 4. Para regenerar as tabelas, reexecute o notebook após modificações nos dados CSV
%% ============================================================================
